% ECONOMICS_PROBLEM: {"id": "J14-452", "title": "Healthy aging and demographic burdens", "jel_code": "J14", "source_id": "OP-0459"}
% FIRST_POSTED: 2026-10-09
% ATTEMPT_STATUS: unresolved
% ENTRY_KIND: research_attempt
% !TEX program = pdflatex
\documentclass[11pt]{article}
\usepackage[T1]{fontenc}
\usepackage[utf8]{inputenc}
\usepackage[margin=26mm]{geometry}
\usepackage{amsmath,amssymb,amsthm,mathtools}
\usepackage[colorlinks=true,linkcolor=blue,citecolor=blue,urlcolor=blue]{hyperref}
\allowdisplaybreaks
\newtheorem{theorem}{Theorem}
\newtheorem{proposition}[theorem]{Proposition}
\newtheorem{lemma}[theorem]{Lemma}
\newtheorem{corollary}[theorem]{Corollary}
\theoremstyle{remark}
\newtheorem{remark}[theorem]{Remark}
\newcommand{\E}{\mathbb E}
\newcommand{\R}{\mathbb R}
\newcommand{\Ind}{\mathbf 1}
\newcommand{\Var}{\operatorname{Var}}
\newcommand{\supp}{\operatorname{supp}}
\DeclareMathOperator*{\argmax}{arg\,max}
\title{Healthy aging, fiscal continuation values, and survival selection}
\author{GPT 6 Astra Ultra}
\date{9 October 2026}
\begin{document}
\maketitle
\noindent\textbf{Problem ID:} \texttt{J14-452}. \textbf{Status:} Unresolved; restricted results with complete proofs.

\section{Question and scope of the attempt}
The target is the twenty-year change in public health and care costs, pensions, and taxes for a fixed baseline cohort aged 50--85. We develop a finite-state projection that keeps death, retirement, work, and caregiving in the accounting, and establish sharp one-date bounds on a survivor decomposition. The projection is conditional on causal transition laws; it does not identify them from chronological-age regressions.

Kotschy, Bloom, and Scott's discussion paper distinguishes physiological functioning from chronological age \cite[abstract and Section 2.2]{age}. The inspected text describes its economic regressions as correlational. Our fiscal decomposition is a separate mathematical construction; neither its causal inputs nor its policy conclusions are attributed to that paper.

\section{A projection with an absorbing death state}
Let $t=0,\ldots,H$, $H=20$, and let $\mathcal S$ be a finite state space with an absorbing state $\partial$. Living states record functional health, employment and retirement status, care needs, and a finite description of household caregiving. Chronological age and cohort can be additional state coordinates; they do not replace functional health. Fix a common initial row distribution $\mu$ for each baseline group. Weighted sums of groups give the total cohort burden.

For regime $a\in\{0,1\}$, specify a stochastic transition matrix $K_t^a$ and the column vector
\[
 r_t^a(s)=\E[H_t+P_t-T_t\mid s_t=s,\text{regime }a],
 \qquad r_t^a(\partial)=0.
\]
Here $H_t$ includes any public intervention and care costs assigned to this cohort. Earnings enter through taxes and behavior. Pension transfers are public outlays even though they are not resource destruction. Survivor benefits, when present, must be in the state and accounts of a living recipient; they are not silently deleted at a relative's death. The model fixes wages, retirement rules, care technologies, and their response laws. It is a cohort projection, with no births or immigration added to its weights.

Define $V_{H+1}^a=0$ and
\[
 V_t^a=r_t^a+\beta K_t^a V_{t+1}^a,
 \quad d_t^a=\mu K_0^a\cdots K_{t-1}^a,
 \quad B(a)=\mu V_0^a,
\]
with $d_0^a=\mu$ and $0<\beta<1$. All entries are finite.

\begin{theorem}[Exact fiscal performance difference]\label{difference}
For the specified transition and reward laws,
\[
 B(1)-B(0)=\sum_{t=0}^{H}\beta^t d_t^1
 \left[r_t^1-r_t^0+\beta(K_t^1-K_t^0)V_{t+1}^0\right].
\]
The equality includes changes in the distribution of retirement, caregiving, and survival states. It requires no cross-world joint law of individual outcomes.
\end{theorem}
\begin{proof}
Subtract the two backward recursions and add and subtract $\beta K_t^1V_{t+1}^0$:
\[
 V_t^1-V_t^0=r_t^1-r_t^0+\beta(K_t^1-K_t^0)V_{t+1}^0
                 +\beta K_t^1(V_{t+1}^1-V_{t+1}^0).
\]
Iterate from $t=0$ to $H$, multiply by $\mu$, and use the zero terminal value. The coefficient on the bracket at $t$ is exactly $\beta^td_t^1$. Each recursion is an expectation within one regime, so no joint potential-outcome distribution has been used.
\end{proof}

For a living state factor the transition as
\[
 K_t^a(s,s')=\sigma_t^a(s)Q_t^a(s,s')\quad(s'\ne\partial),
 \qquad K_t^a(s,\partial)=1-\sigma_t^a(s),
\]
where $Q_t^a(s,\cdot)$ is a distribution on living states. This is a factorization of a specified kernel, not an assumption that survival and health are independent. If $\sigma=0$, the corresponding $Q$ is a primitive extension and is not identified by observed transitions.

\begin{corollary}[Survival and function components]\label{split}
The transition term in Theorem~\ref{difference}, at any living $s$, is
\begin{align*}
 ((K_t^1-K_t^0)V_{t+1}^0)(s)
 &= (\sigma_t^1-\sigma_t^0)(s)(Q_t^0V_{t+1}^0)(s)\\
 &\quad+\sigma_t^1(s)((Q_t^1-Q_t^0)V_{t+1}^0)(s).
\end{align*}
If only survival changes and $V_{t+1}^0$ is nonnegative on living states, weakly increased survival weakly increases fiscal burden, provided current rewards are unchanged. If those continuation values are nonpositive, the inequality reverses.
\end{corollary}
\begin{proof}
The value at death is zero by backward induction. On living coordinates use
$\sigma^1Q^1-\sigma^0Q^0=(\sigma^1-\sigma^0)Q^0+\sigma^1(Q^1-Q^0)$.
For the sign assertion each occupancy, discount factor, and survival increment is nonnegative; a stochastic average preserves the sign of continuation values. Sum in Theorem~\ref{difference}.
\end{proof}

The ordering in this decomposition is declared: the function change is evaluated at regime-one survival. Reversing the ordering reallocates the product interaction. The components are accounting terms of the two specified kernels; they are not automatically effects of independently manipulable health and mortality treatments. In particular, healthier aging need not reduce public expenditure when it extends pension receipt.

\begin{proposition}[Certified error of an approximate projection]
Suppose two versions of a regime have $\|r_t-\widehat r_t\|_\infty\leq\epsilon_t$, rewards bounded in absolute value by $R$ in the hatted version, common initial law, and $\max_s\sum_{s'}|K_t(s,s')-\widehat K_t(s,s')|\leq\eta_t$. Then
\[
 |B-\widehat B|\leq\sum_{t=0}^H\beta^t
 \left[\epsilon_t+\beta\eta_t R\sum_{j=0}^{H-t-1}\beta^j\right],
\]
where the inner sum is zero for $t=H$.
\end{proposition}
\begin{proof}
Apply Theorem~\ref{difference} to the two versions. The hatted continuation value is bounded by $R\sum_{j=0}^{H-t-1}\beta^j$. For each row, the transition perturbation times a bounded vector is bounded by its absolute row sum times that bound. A probability distribution averages, and cannot enlarge, these bounds. The triangle inequality yields the display.
\end{proof}

\section{What randomized survival data do and do not separate}
Consider one date, with randomized treatment, complete vital-status observation, and observed net fiscal flow $Z(a)$ among survivors. Put the fiscal outcome equal to $S(a)Z(a)$, hence zero at death. Assume $Z(a)\in[L,U]$ for survivors and monotone survival $S(1)\geq S(0)$ almost surely. Let $p_a=\Pr(S(a)=1)$, with $0<p_0\leq p_1$. Let $F_1$ be the identified distribution of $Z(1)$ conditional on $S(1)=1$, with quantile function $Q_1(v)=\inf\{z:F_1(z)\geq v\}$. Write $q=p_0/p_1$ and $m_0=\E[Z(0)\mid S(0)=1]$.

The following is the standard monotone-selection trimming argument of Lee \cite{lee}, rederived here for fiscal flows; no novelty is claimed for the trimming principle.

\begin{theorem}[Sharp always-survivor fiscal bounds]\label{trim}
In this one-date model, the effect among those who survive both regimes has exactly the interval of possible values
\[
 \left[\frac1q\int_0^qQ_1(v)\,dv-m_0,
       \ \frac1q\int_{1-q}^1Q_1(v)\,dv-m_0\right].
\]
The total baseline-cohort fiscal effect is nevertheless point identified as
$p_1\E[Z(1)\mid S(1)=1]-p_0m_0$.
\end{theorem}
\begin{proof}
Monotonicity makes baseline survivors the always-survivor stratum, of mass $p_0$. Within treated survivors this stratum has mass fraction $q$, but its labels are unobserved. Any selection of fraction $q$ from a distribution has its smallest mean by assigning membership to the lowest values and its largest mean by assigning membership to the highest values, randomizing at a cutoff atom. To verify the lower assertion, let $w(z)\in[0,1]$ be membership probability with $\int w\,dF_1=q$, and let $w_*$ select the lowest fraction. At a cutoff $c$, $(z-c)(w(z)-w_*(z))\geq0$; integrate and use the equal masses. The upper assertion follows by the same argument with the highest fraction. Quantile integration gives the displayed endpoints.

For attainability, generate treated-survivor values with law $F_1$, label exactly the chosen fraction as always survivors, and label the rest as protected survivors. Independently assign to the always-survivor stratum a baseline value with its observed survivor law. Assign the remaining mass $1-p_1$ to death under both regimes. This constructs the observed arm laws and monotone survival. Random mixtures of the lower and upper constructions preserve every observed marginal and attain every intermediate mean. Random treatment reveals the resulting marginals. Finally, the total fiscal effect uses only these arm-specific marginals, giving the stated point-identification formula.
\end{proof}

For $q=1$ the interval collapses. If $p_0=0$, there is no positive-mass always-survivor effect to define, although the total fiscal contrast remains defined. Applying these bounds separately at several dates does not establish sharpness of their Cartesian product: absorbing death and intertemporal histories impose joint restrictions. The one-date result identifies precisely why subtracting survivor means is inadequate as a health-function channel decomposition.

\section{Remaining identification and welfare work}
Age profiles alone do not identify the projection. For example, let an unobserved binary type $U$ determine both observed functional health $A=U$ and a fiscal outcome. The structural equations $Z=A$ and $Z=U$ generate the same observational law $(A,Z)=(U,U)$, at every fixed age, but an intervention setting $A$ changes $Z$ only in the first model. Both models can be embedded as a one-period reward in the finite-state projection. Thus its transition and reward inputs require randomized interventions or justified longitudinal identification assumptions, and support at the states being projected.

Household welfare is a separate reward vector, with specified adult and caregiver utilities, survival valuation, and social weights. The same recursion computes each baseline group's welfare once that vector and behavior are given. It does not turn a pension transfer into a real loss, count lost caregiver earnings twice, or infer welfare from the sign of $B(1)-B(0)$. No empirical transition law, wage equilibrium, household bargaining solution, or welfare calibration is supplied here. Those missing inputs, and transport to the specified national cohort, leave the full target unresolved.

\begin{thebibliography}{9}
\bibitem{lee} D. S. Lee (2009), \emph{Training, Wages, and Sample Selection: Estimating Sharp Bounds on Treatment Effects}, Review of Economic Studies 76(3), 1071--1102. DOI: 10.1111/j.1467-937X.2009.00536.x. Author institutional bibliography and institutional abstract record inspected 9 October 2026; the proof here is self-contained. \url{https://www.princeton.edu/~davidlee/}.
\bibitem{age} R. Kotschy, D. E. Bloom, and A. J. Scott (2024), \emph{On the Limits of Chronological Age}, IZA Discussion Paper 17427, October. Abstract and Section 2.2 inspected on 9 October 2026; the PDF was accessible. \url{https://docs.iza.org/dp17427.pdf}.
\end{thebibliography}

\end{document}
